\chapter{定積分の対称性}\label{章:定積分の対称性}

不定積分では，微分すると与えられた関数になる式を探した。定積分でも原始関数が見つかれば計算できるが，定積分にはもう一つ情報がある。どこからどこまで積分するかという区間である。たとえば，同じ関数でも区間が原点の左右に同じ長さだけ延びていれば，左右の値を組にして考えられる。区間の端が互いに逆数なら，逆数の置換によって両端を交換できる。

この章では，置換後の積分をすぐに計算せず，置換前の積分と並べてみる。二つを加えたとき，分母が消えることもあれば，変数を掛けた部分が定数に変わることもある。そのため，原始関数を求めるには手間がかかる式から，短い計算で定積分の値が決まる。

以下では，特に断らない限り，被積分関数は積分する閉区間で連続であるとする。途中に分母が零になる点がある式に，対称性だけを見て同じ議論を当てはめることはできない。特異点を含む積分は，後の章で定義を確かめてから扱う。

\section{左右の値を組にする}\label{節:奇偶の対称性}

\begin{definition}{偶関数と奇関数}
原点に関して対称な範囲で定義された関数 $f$ について，すべての $x$ で
\[
f(-x)=f(x)
\]
が成り立つとき，$f$ を偶関数という。また，すべての $x$ で
\[
f(-x)=-f(x)
\]
が成り立つとき，$f$ を奇関数という。
\end{definition}

偶関数のグラフは $y$ 軸について対称であり，奇関数のグラフは原点について対称である。$x^2$ と $\cos x$ は偶関数，$x^3$ と $\sin x$ は奇関数である。ただし，式の中に $\sin$ があれば奇関数になるわけではない。$\sin(x^2)$ では $x$ を $-x$ に変えても $x^2$ が変わらないので偶関数である。式全体に $-x$ を代入することが，確実な判定になる。

$a>0$ とし，$[-a,a]$ の積分を零で分ける。左半分で $x=-t$ と置くと，$x=-a,0$ に対して $t=a,0$ であり，$\dd x=-\dd t$ だから，
\[
\int_{-a}^{0}f(x)\dd x
 =\int_a^0 f(-t)(-\dd t)
 =\int_0^a f(-t)\dd t.
\]
積分変数の文字をそろえれば，
\begin{equation}\label{式:左右の積分}
\int_{-a}^a f(x)\dd x
 =\int_0^a\{f(x)+f(-x)\}\dd x
\end{equation}
となる。これは偶関数でも奇関数でもない場合にも使える式である。偶関数なら中括弧が $2f(x)$ になり，奇関数なら零になる。

\begin{theorem}{対称区間での偶関数・奇関数の積分}
$a>0$ とする。$f$ が偶関数なら
\[
\int_{-a}^a f(x)\dd x=2\int_0^a f(x)\dd x,
\]
$f$ が奇関数なら
\[
\int_{-a}^a f(x)\dd x=0
\]
である。
\end{theorem}

図\ref{図:奇偶の積分}の右図では，左右の部分の幾何学的な面積は等しい。しかし左側の定積分は負，右側は正である。奇関数の積分が零という結論は，図形の面積が零という意味ではない。面積を求めるなら絶対値を積分する必要がある。

\begin{figure}[htbp]
\centering
\begin{tikzpicture}[scale=1.05,>=Stealth]
\begin{scope}[xshift=-3.1cm]
\fill[blue!12] plot[domain=-1.35:1.35,samples=50](\x,{0.65*\x*\x})--(1.35,0)--(-1.35,0)--cycle;
\draw[->](-1.7,0)--(1.8,0)node[right]{$x$};
\draw[->](0,-0.2)--(0,1.7)node[above]{$y$};
\draw[blue!70!black,thick] plot[domain=-1.45:1.45,samples=60](\x,{0.65*\x*\x});
\draw[dashed](-1.35,0)--(-1.35,1.185);
\draw[dashed](1.35,0)--(1.35,1.185);
\node[below]at(-1.35,0){$-a$};\node[below]at(1.35,0){$a$};
\node at(0,-0.65){偶関数：左右を加える};
\end{scope}
\begin{scope}[xshift=1.8cm]
\fill[blue!12](0,0)--plot[domain=0:1.35,samples=40](\x,{0.55*\x*\x*\x})--(1.35,0)--cycle;
\fill[orange!18](0,0)--plot[domain=0:-1.35,samples=40](\x,{0.55*\x*\x*\x})--(-1.35,0)--cycle;
\draw[->](-1.7,0)--(1.8,0)node[right]{$x$};
\draw[->](0,-1.55)--(0,1.7)node[above]{$y$};
\draw[blue!70!black,thick] plot[domain=-1.4:1.4,samples=60](\x,{0.55*\x*\x*\x});
\draw[dashed](-1.35,0)--(-1.35,-1.353);
\draw[dashed](1.35,0)--(1.35,1.353);
\node[above]at(-1.35,0){$-a$};\node[below]at(1.35,0){$a$};
\node at(0,-1.95){奇関数：符号付きの量が打ち消す};
\end{scope}
\end{tikzpicture}
\caption{対称性と，定積分の符号の関係。曲線はそれぞれ $x^2$ 型，$x^3$ 型である。}\label{図:奇偶の積分}
\end{figure}

\begin{example}{混じった式から奇関数の部分を除く}
次の定積分を求めよ。
\[
I=\int_{-2}^{2}\left(x^3\e^{-x^2}+2x^2+1\right)\dd x.
\]
\end{example}

$x^3\e^{-x^2}$ の原始関数を先に探すこともできるが，区間が $[-2,2]$ なので，まず $x$ を $-x$ に変えてみる。$x^3$ の符号だけが変わり，$\e^{-x^2}$ は変わらない。したがって積 $x^3\e^{-x^2}$ は奇関数で，その積分は零である。残った $2x^2+1$ は偶関数だから，
\[
I=2\int_0^2(2x^2+1)\dd x
 =2\left[\frac{2x^3}{3}+x\right]_0^2
 =\frac{44}{3}.
\]
ここでは式全体が偶関数か奇関数かを決める必要はなかった。和を分けて，対称性が分かる部分にだけ使えばよい。

さらに，一般の $f$ に対して
\[
f_{\mathrm e}(x)=\frac{f(x)+f(-x)}{2},\qquad
f_{\mathrm o}(x)=\frac{f(x)-f(-x)}{2}
\]
と置くと，$f=f_{\mathrm e}+f_{\mathrm o}$ であり，前者は偶関数，後者は奇関数になる。よって対称区間の定積分に寄与するのは $f_{\mathrm e}$ の部分だけである。これは新しい積分公式を覚えるための分解ではなく，式\eqref{式:左右の積分}が何を選び出しているかを示す分解である。

\section{区間の両端を交換する}\label{節:区間の反転}

対称の中心は原点に限らない。区間 $[a,b]$ の両端を交換するには，$x=a+b-t$ と置けばよい。対応は
\[
x=a\longleftrightarrow t=b,\qquad
x=b\longleftrightarrow t=a
\]
であり，$\dd x=-\dd t$ だから，
\begin{equation}\label{式:区間反転}
\int_a^b f(x)\dd x=\int_a^b f(a+b-x)\dd x
\end{equation}
を得る。置換の直後は積分の上下端が逆になるが，微分の負号と合わせて向きを戻している。

この式の左右は同じ値なので，加えて二で割ると
\begin{equation}\label{式:反転平均}
\int_a^b f(x)\dd x
 =\frac12\int_a^b\{f(x)+f(a+b-x)\}\dd x
\end{equation}
となる。$f(x)+f(a+b-x)$ が定数 $C$ になれば，答えは $C(b-a)$ の半分である。元の式を個別に積分できるかどうかは，この計算では問題にならない。

今度は，関数 $g$ が区間の中央について対称，すなわち
\[
g(a+b-x)=g(x)
\]
を満たすとする。このとき $xg(x)$ の積分を $I$ とすると，反転により
\[
I=\int_a^b(a+b-x)g(x)\dd x.
\]
元の式を加えると $x+(a+b-x)=a+b$ なので，
\begin{equation}\label{式:対称な重みの一次因子}
\int_a^b xg(x)\dd x
 =\frac{a+b}{2}\int_a^b g(x)\dd x
\end{equation}
となる。中央からの距離が等しい二点では $g$ の値が等しく，その二点の $x$ 座標の平均は常に区間の中央になる，という関係である。

\begin{example}{変数が掛かった三角関数}
次の定積分を求めよ。
\[
I=\int_0^\pi x\sin^3x\dd x.
\]
\end{example}

$x$ と三角関数の積を見ると部分積分を考えたくなる。しかしこの区間では $\sin(\pi-x)=\sin x$ である。$g(x)=\sin^3x$ は区間の中央について対称なので，式\eqref{式:対称な重みの一次因子}から
\[
I=\frac{\pi}{2}\int_0^\pi\sin^3x\dd x.
\]
残りは $\sin^3x=(1-\cos^2x)\sin x$ として $u=\cos x$ と置く。
\[
\int_0^\pi\sin^3x\dd x
 =\int_{-1}^1(1-u^2)\dd u
 =\left[u-\frac{u^3}{3}\right]_{-1}^1
 =\frac43.
\]
したがって $I=\frac{2\pi}{3}$ である。反転で $x$ を取り除き，置換で三角関数を取り除くという二段階で計算できた。

式\eqref{式:対称な重みの一次因子}を使う際には，$xg(x)$ 全体が対称なのではなく，$g$ が対称であることを確認する。たとえば $\int_0^1 x\e^x\dd x$ では $\e^{1-x}=\e^x$ が成り立たないので，$x$ を区間の中央 $\frac12$ に置き換えることはできない。

また，$x^2g(x)$ まで同じように中央の二乗に置き換えられるわけでもない。$x^2+(a+b-x)^2$ は $x$ を含むからである。反転して足すことで何が消えるのかを，その都度式で確かめる必要がある。

\section{正弦と余弦を交換する}\label{節:正弦余弦の交換}

区間 $\left[0,\frac{\pi}{2}\right]$ での反転は，$x$ を $\frac{\pi}{2}-x$ に変える操作である。このとき
\[
\sin\left(\frac{\pi}{2}-x\right)=\cos x,\qquad
\cos\left(\frac{\pi}{2}-x\right)=\sin x
\]
なので，正弦と余弦が入れ替わる。分子と分母に両者が現れる場合，入れ替えた式との和を調べるとよい。

\begin{example}{分母を変えずに分子を交換する}
次の定積分を求めよ。
\[
I=\int_0^{\frac{\pi}{2}}
\frac{\sin^4x}{\sin^4x+\cos^4x}\dd x.
\]
\end{example}

$\sin x$ と $\cos x$ は同時に零にならないので，分母は区間全体で正である。反転すると分子は $\cos^4x$ に変わり，分母は変わらない。したがって
\[
I=\int_0^{\frac{\pi}{2}}
\frac{\cos^4x}{\sin^4x+\cos^4x}\dd x.
\]
二つを加えれば，分子の和が分母そのものになる。
\[
2I=\int_0^{\frac{\pi}{2}}1\dd x=\frac{\pi}{2},
\qquad I=\frac{\pi}{4}.
\]
$\sin^4x$ の次数を下げたり，分母を展開したりする前に，分子の足し合わせで分母ができることに気づくと計算を減らせる。

この関係は四乗に特有ではない。$q$ を $[0,1]$ 上の連続関数とし，$q(\sin x)+q(\cos x)$ が区間内で零にならなければ，
\[
\int_0^{\frac{\pi}{2}}
\frac{q(\sin x)}{q(\sin x)+q(\cos x)}\dd x
 =\frac{\pi}{4}
\]
である。たとえば $q(t)=t^2+1$ でも同じ値になる。重要なのは指数の値ではなく，反転で分子の二つの候補が交換されることである。

図\ref{図:相補重み}は，例題の被積分関数を $w(x)$ としたグラフである。左右を反転したグラフ $w(\frac{\pi}{2}-x)$ と，同じ $x$ における高さを足すと一になる。それぞれの積分が等しいので，全体の長方形の面積を二等分している。

\begin{figure}[htbp]
\centering
\begin{tikzpicture}[x=3.5cm,y=2.3cm,>=Stealth]
\fill[blue!10](0,0)--plot[domain=0:1.570796,samples=100](\x,{pow(sin(\x r),4)/(pow(sin(\x r),4)+pow(cos(\x r),4))})--(1.570796,0)--cycle;
\draw[->](-0.08,0)--(1.73,0)node[right]{$x$};
\draw[->](0,-0.08)--(0,1.2)node[above]{$y$};
\draw[dashed](0,1)--(1.570796,1)--(1.570796,0);
\draw[dashed](0.785398,0)--(0.785398,1);
\draw[blue!75!black,thick]plot[domain=0:1.570796,samples=100](\x,{pow(sin(\x r),4)/(pow(sin(\x r),4)+pow(cos(\x r),4))});
\draw[orange!80!black,thick,dashed]plot[domain=0:1.570796,samples=100](\x,{pow(cos(\x r),4)/(pow(sin(\x r),4)+pow(cos(\x r),4))});
\node[left]at(0,1){$1$};\node[below]at(0,0){$0$};
\node[below]at(0.785398,0){$\frac{\pi}{4}$};
\node[below]at(1.570796,0){$\frac{\pi}{2}$};
\node[blue!75!black]at(1.15,0.65){$w(x)$};
\node[orange!80!black]at(0.48,1.26){$w(\frac{\pi}{2}-x)$};
\end{tikzpicture}
\caption{反転した二つの重みは，同じ場所で足すと一になる。}\label{図:相補重み}
\end{figure}

\section{周期があるとき，どこから一周期を取るか}\label{節:周期関数の積分}

$T>0$ に対して $f(x+T)=f(x)$ がすべての $x$ で成り立つとき，$T$ を $f$ の周期という。周期という言葉だけから，$T$ は最小の正の周期だとは限らない。$\sin x$ の周期として $2\pi$ も $4\pi$ も使える。

周期関数をちょうど一周期分だけ積分した値は，始点によらない。まず $0\leq a\leq T$ の場合，
\[
\int_a^{a+T}f(x)\dd x
 =\int_a^T f(x)\dd x+\int_T^{a+T}f(x)\dd x.
\]
第二項で $x=t+T$ と置くと $f(t+T)=f(t)$ だから，
\[
\int_T^{a+T}f(x)\dd x=\int_0^a f(t)\dd t.
\]
二つを足すと
\begin{equation}\label{式:一周期の積分}
\int_a^{a+T}f(x)\dd x=\int_0^T f(x)\dd x
\end{equation}
になる。一般の $a$ も，整数個の $T$ だけ平行移動すれば始点を $[0,T)$ に移せるため，同じ結論になる。区間を長さ $T$ の部分に分ければ，正の整数 $n$ に対して
\[
\int_a^{a+nT}f(x)\dd x=n\int_0^T f(x)\dd x
\]
も得られる。

周期があるというだけでは，一周期分の積分が零とは限らない。たとえば $\sin^2x$ は周期 $\pi$ を持つが，一周期の積分は $\frac{\pi}{2}$ である。符号付きの部分が打ち消し合うか，積分区間が周期の整数倍になっているかを別々に確認する。

\begin{example}{絶対値がある周期関数}
次の定積分を求めよ。
\[
I=\int_{-\frac{\pi}{6}}^{\frac{13\pi}{6}}|\sin x|\dd x.
\]
\end{example}

$\sin(x+\pi)=-\sin x$ なので，絶対値を取ると周期は $\pi$ になる。一周期分は
\[
\int_0^\pi|\sin x|\dd x=\int_0^\pi\sin x\dd x=2.
\]
与えられた区間の長さは $\frac{7\pi}{3}$ であり，$2\pi$ と $\frac{\pi}{3}$ に分けられる。まず $-\frac{\pi}{6}$ から $\frac{11\pi}{6}$ までは二周期分なので積分は $4$ である。残りでは $x=t+2\pi$ と置き，
\[
\int_{\frac{11\pi}{6}}^{\frac{13\pi}{6}}|\sin x|\dd x
 =\int_{-\frac{\pi}{6}}^{\frac{\pi}{6}}|\sin t|\dd t.
\]
$|\sin t|$ は偶関数で，$0\leq t\leq\frac{\pi}{6}$ では $\sin t\geq0$ だから，
\[
I=4+2\int_0^{\frac{\pi}{6}}\sin t\dd t
 =4+2\left(1-\frac{\sqrt3}{2}\right)
 =6-\sqrt3.
\]
周期によって長い区間を短くし，偶奇によって符号の変わる区間を半分にした。周期の整数倍に満たない残りは，始点によって値が変わるので省いてはいけない。

\section{逆数の置換で端を交換する}\label{節:逆数対称性}

区間の端が $\frac1b$ と $b$ であるときは，$x=\frac1t$ でも両端を交換できる。ただし反転 $x=a+b-t$ と異なり，微分は
\[
\dd x=-\frac{1}{t^2}\dd t
\]
になる。関数の $x$ を逆数に変えるだけでは，正しい置換にならない。

特に積分が $\frac{\dd x}{x}$ を含むと，
\[
\frac{\dd x}{x}=-\frac{\dd t}{t}
\]
となり，同じ形に戻る。$b>1$ とすると，
\begin{equation}\label{式:逆数置換の対称性}
\int_{\frac1b}^b F(x)\frac{\dd x}{x}
 =\int_{\frac1b}^b F\left(\frac1x\right)\frac{\dd x}{x}
\end{equation}
である。したがって，$F(x)+F(\frac1x)=C$ なら
\[
2\int_{\frac1b}^b F(x)\frac{\dd x}{x}
 =C\int_{\frac1b}^b\frac{\dd x}{x}
 =2C\log b
\]
となる。

\begin{example}{逆数にすると補い合う分数}
次の定積分を求めよ。
\[
I=\int_{\frac13}^3\frac{1}{x(1+x^2)}\dd x.
\]
\end{example}

区間の両端が互いに逆数であり，分母に $x$ がある。そこで $F(x)=\frac{1}{1+x^2}$ と見て，$\frac{\dd x}{x}$ を残す。
\[
F\left(\frac1x\right)=\cfrac{1}{1+\cfrac{1}{x^2}}
 =\frac{x^2}{1+x^2},
\qquad F(x)+F\left(\frac1x\right)=1.
\]
式\eqref{式:逆数置換の対称性}を使えば
\[
2I=\int_{\frac13}^3\frac{\dd x}{x}
 =\left[\log x\right]_{\frac13}^3=2\log3,
\]
よって $I=\log3$ である。

この例を直接置換しても，
\[
I=\int_3^{\frac13}
\cfrac{1}{\cfrac1t(1+\cfrac1{t^2})}
\left(-\frac{1}{t^2}\right)\dd t
=\int_{\frac13}^3\frac{t}{1+t^2}\dd t
\]
となる。元の被積分関数との和は $\frac1x$ であり，上の計算と一致する。微分因子を忘れると，この和は現れない。

逆数の対称性を，いつもの左右の対称性として見直すこともできる。$x=\e^u$ と置くと，$\frac{\dd x}{x}=\dd u$，区間は $[-\log b,\log b]$ になる。さらに $x$ を $\frac1x$ に変える操作は，$u$ を $-u$ に変える操作になる。数直線上の普通の長さでは $\frac1b$ と $b$ は $1$ から等距離ではないが，対数で測った座標では左右対称なのである。

一般に，端が正で $ab=k$ なら，$x=\frac{k}{t}$ は $[a,b]$ の両端を交換する。この場合にも $\frac{\dd x}{x}=-\frac{\dd t}{t}$ が成り立つ。端を見て置換を考えるときは，和が一定か，積が一定かが手がかりになる。

\section{足すと一になる重み}\label{節:相補重み}

次の二つの式を比べる。
\[
\frac{1}{1+\e^x},\qquad
\frac{1}{1+\e^{-x}}=\frac{\e^x}{1+\e^x}.
\]
分子の和が分母なので，二つを足すと一になる。そこで
\[
w(x)=\frac{1}{1+\e^x}
\]
と置けば $w(x)+w(-x)=1$ である。このように，対称な二点で値の和が一になる因子を，ここでは互いに補い合う重みとして考える。

$g$ が偶関数なら，
\[
I=\int_{-a}^a g(x)w(x)\dd x
 =\int_{-a}^a g(x)w(-x)\dd x
\]
である。二式を加えると
\begin{equation}\label{式:偶関数と相補重み}
\int_{-a}^a\frac{g(x)}{1+\e^x}\dd x
 =\frac12\int_{-a}^a g(x)\dd x
 =\int_0^a g(x)\dd x
\end{equation}
となる。指数関数の因子を一度も積分せずに，式から取り除けた。

さらに，$h$ が奇関数なら $h(-x)=-h(x)$ なので，$\frac{1}{1+\e^{h(x)}}$ にも同じ議論が使える。必要なのは指数が $x$ であることではなく，反転すると指数の符号が変わることである。

\begin{example}{指数が複雑でも，反転で符号が変わる}
次の定積分を求めよ。
\[
I=\int_{-2}^2\frac{x^4+\cos x}{1+\e^{x^3+x}}\dd x.
\]
\end{example}

分子 $g(x)=x^4+\cos x$ は偶関数で，指数 $h(x)=x^3+x$ は奇関数である。$x$ を $-x$ に変えると分子はそのまま，指数だけが負になる。したがって
\[
2I=\int_{-2}^2(x^4+\cos x)\dd x,
\]
よって
\[
I=\int_0^2(x^4+\cos x)\dd x
 =\left[\frac{x^5}{5}+\sin x\right]_0^2
 =\frac{32}{5}+\sin2.
\]
分子は区間全体で正である必要はない。偶関数であることだけが，上の等式に使われている。

一方，分子が奇関数のときは，同じ結論にはならない。たとえば
\[
K=\int_{-a}^a\frac{x}{1+\e^x}\dd x
\]
を左右で組にすると，
\[
K=\int_0^a x\left(\frac{1}{1+\e^x}-\frac{1}{1+\e^{-x}}\right)\dd x
 =\int_0^a\frac{x(1-\e^x)}{1+\e^x}\dd x.
\]
$a>0$ なら被積分関数は $0<x\leq a$ で負なので，$K<0$ である。「分子の奇関数を対称区間で積分するから零」と考えると，分母が偶関数でないことを見落とす。積全体の対称性か，反転後の式との関係を確かめる必要がある。

この考えは原点以外にも移せる。区間の中央を $c=\frac{a+b}{2}$ とし，$g(a+b-x)=g(x)$，$w(x)+w(a+b-x)=1$ なら
\[
\int_a^b g(x)w(x)\dd x=\frac12\int_a^b g(x)\dd x.
\]
たとえば $w(x)=\frac{1}{1+\e^{x-c}}$ はこの条件を満たす。変数を $u=x-c$ に移せば，原点に関する場合と全く同じ証明になる。

\section{対称性で何を簡単にするか}\label{節:定積分の手法選択}

ここまでの方法は，どれも新しい原始関数を増やすものではない。定積分を，値の等しい別の積分に書き換え，二つを組にする方法である。最初に区間を見て，次にその区間を保つ変換を考えると，計算の候補を絞れる。

\begin{center}
\begin{tabularx}{0.96\linewidth}{@{}lXX@{}}
\toprule
区間や式の特徴 & 試す変換 & 比べる部分\\
\midrule
$[-a,a]$ & $x\mapsto -x$ & 偶関数と奇関数の部分\\
$[a,b]$ & $x\mapsto a+b-x$ & 和が定数になるか，一次因子が消えるか\\
$\left[0,\frac{\pi}{2}\right]$ & $x\mapsto\frac{\pi}{2}-x$ & 正弦と余弦が交換されるか\\
長い区間と周期関数 & 整数周期だけ平行移動 & 完全な周期と残りの区間\\
$\left[\frac1b,b\right]$ と $\frac{\dd x}{x}$ & $x\mapsto\frac1x$ & 逆数にした関数との和\\
\bottomrule
\end{tabularx}
\end{center}

候補が見つかっても，対称性だけで最後まで解けるとは限らない。対称性によって厄介な因子を一つ除き，残った部分に置換積分や部分積分を使えばよい。

\begin{example}{反転の後で置換と部分分数を使う}
次の定積分を求めよ。
\[
I=\int_0^\pi\frac{x\sin x}{1+\sin^2x}\dd x.
\]
\end{example}

$\frac{\sin x}{1+\sin^2x}$ は $x$ を $\pi-x$ に変えても同じである。したがって，最初に $x$ を対称性で取り除く。
\[
I=\frac{\pi}{2}J,\qquad
J=\int_0^\pi\frac{\sin x}{1+\sin^2x}\dd x.
\]
$J$ では分子の $\sin x\dd x$ が $-\dd(\cos x)$ に対応する。$u=\cos x$ と置くと，$1+\sin^2x=2-u^2$ だから
\[
J=\int_{-1}^1\frac{1}{2-u^2}\dd u.
\]
この分母は $(\sqrt2-u)(\sqrt2+u)$ に分かれる。よって
\[
\frac{1}{2-u^2}
 =\frac{1}{2\sqrt2}\left(\frac{1}{\sqrt2-u}+\frac{1}{\sqrt2+u}\right),
\]
\[
J=\left[\frac{1}{2\sqrt2}
\log\frac{\sqrt2+u}{\sqrt2-u}\right]_{-1}^1
 =\frac{1}{\sqrt2}\log\frac{\sqrt2+1}{\sqrt2-1}
 =\sqrt2\log(\sqrt2+1).
\]
ここでは $-1\leq u\leq1$ で対数の分子・分母がともに正である。最後の等式では $(\sqrt2+1)(\sqrt2-1)=1$ を使った。したがって
\[
I=\frac{\pi}{\sqrt2}\log(\sqrt2+1).
\]

もし最初に部分積分を使うなら，$\frac{\sin x}{1+\sin^2x}$ の原始関数が必要になり，さらにその対数を積分する項が現れる。反転によって $x$ を除いたため，必要な計算が一度の置換と部分分数で済んだ。同じ式でも，不定積分では「微分の逆を作れるか」を見て，定積分ではそれに加えて「区間を保つ置換で二つの式を補い合わせられるか」を見るのである。

\section*{章末問題}
\begin{enumerate}[label=\textbf{\thechapter.\arabic*},leftmargin=*]
\item\label{問:奇偶の混合}
次の定積分を求めよ。
\[
\int_{-2}^2\left(\frac{x^3}{1+x^2}+2x^2-\cos x\right)\dd x.
\]
\item\label{問:一次重み反転}
反転を利用して，$\int_0^\pi x\sin^2x\dd x$ を求めよ。
\item\label{問:三角相補}
$\int_0^{\frac{\pi}{2}}\frac{\cos^6x}{\sin^6x+\cos^6x}\dd x$ を求めよ。分母が零にならないことも確かめよ。
\item\label{問:周期絶対値}
$\int_{\frac{\pi}{3}}^{\frac{10\pi}{3}}|\cos x|\dd x$ と
$\int_{\frac{\pi}{3}}^{\frac{10\pi}{3}}\cos x\dd x$ をそれぞれ求め，値が異なる理由を述べよ。
\item\label{問:逆数四乗}
$\int_{\frac14}^{4}\frac{1}{x(1+x^4)}\dd x$ を，逆数の置換を用いて求めよ。
\item\label{問:指数相補}
$\int_{-1}^1\frac{3x^2+1}{1+\e^{x^5}}\dd x$ を求めよ。
\item\label{問:中心を移す相補}
$\int_2^6\frac{(x-4)^2+2}{1+\e^{x-4}}\dd x$ を求めよ。
\item\label{問:二次重みの限界}
$I=\int_0^\pi x^2\sin x\dd x$ とする。反転によって
$I=\int_0^\pi(\pi-x)^2\sin x\dd x$ と書けるが，この二式を加えるだけで $I$ は決まらないことを説明せよ。その後，部分積分を用いて $I$ を求めよ。
\end{enumerate}

\smallskip\noindent 解答は\pageref{解答:第7章}ページ。
